\unitchapter{2}{2}{Computer System Architecture}{p2-02}

\unitmeta{Expected questions: 8--10 · Target: 8+ · Type: numericals (cache,
pipeline, number systems) + concepts (Morris Mano)}

\begin{sylbox}

\begin{itemize}
\tightlist
\item[$\square$]
  Digital logic circuits \& components: gates, Boolean algebra, K-maps, combinational circuits, flip-flops, sequential circuits, ICs, decoders, multiplexers, registers, counters, memory unit
\item[$\square$]
  Data representation: number systems, complements, fixed \& floating point, error-detection codes, computer arithmetic
\item[$\square$]
  Register transfer \& micro-operations
\item[$\square$]
  Basic computer organization \& design: instruction codes, registers, timing \& control, instruction cycle, interrupts
\item[$\square$]
  Programming the basic computer; microprogrammed control
\item[$\square$]
  CPU: general register \& stack organisation, instruction formats, addressing modes, RISC vs CISC
\item[$\square$]
  Pipeline \& vector processing
\item[$\square$]
  I/O organisation: interfaces, asynchronous transfer, modes of transfer, priority interrupt, DMA, serial communication
\item[$\square$]
  Memory hierarchy: main, auxiliary, associative, cache, virtual memory
\item[$\square$]
  Multiprocessors: interconnection structures, arbitration, synchronisation, cache coherence, multicore
\end{itemize}

\end{sylbox}

\needspace{166pt}

\hypertarget{p2-02--1-digital-logic}{%
\section{Digital Logic}\label{p2-02--1-digital-logic}}

\needspace{130pt}

\hypertarget{p2-02--11-gates}{%
\subsection{Gates}\label{p2-02--11-gates}}

\begingroup\small

\begin{longtable}[]{@{}
  >{\raggedright\arraybackslash}p{(\columnwidth - 4\tabcolsep) * \real{0.0793}}
  >{\raggedright\arraybackslash}p{(\columnwidth - 4\tabcolsep) * \real{0.1874}}
  >{\raggedright\arraybackslash}p{(\columnwidth - 4\tabcolsep) * \real{0.2925}}@{}}
\toprule\noalign{}
\begin{minipage}[b]{\linewidth}\raggedright
\textbf{Gate}
\end{minipage} & \begin{minipage}[b]{\linewidth}\raggedright
\textbf{Expression}
\end{minipage} & \begin{minipage}[b]{\linewidth}\raggedright
\textbf{Output 1 when}
\end{minipage} \\
\midrule\noalign{}
\endhead
\bottomrule\noalign{}
\endlastfoot
AND & AB & all inputs 1 \\
OR & A + B & any input 1 \\
NOT & A\textquotesingle{} & input 0 \\
NAND & (AB)\textquotesingle{} & any input 0 --- \textbf{universal} \\
NOR & (A+B)\textquotesingle{} & all inputs 0 --- \textbf{universal} \\
XOR & A⊕B = A\textquotesingle B + AB\textquotesingle{} & inputs differ (odd number of 1s) \\
XNOR & (A⊕B)\textquotesingle{} & inputs equal \\
\end{longtable}

\endgroup

Gates needed using only NAND: NOT = 1, AND = 2, OR = 3, XOR = \textbf{4}, XNOR = 5.
Using only NOR: NOT = 1, OR = 2, AND = 3, XNOR = \textbf{4}, XOR = 5.

\needspace{137pt}

\hypertarget{p2-02--12-k-map-simplification}{%
\subsection{K-Map Simplification}\label{p2-02--12-k-map-simplification}}

\begin{listingblock}{87pt}
\begin{Verbatim}[fontsize=\fontsize{7.8}{9.5}\selectfont]
 4-variable K-map (Gray code order)          Example: F = Σm(0,2,5,7,8,10,13,15)
          CD                                             CD
 AB    00  01  11  10                          AB    00  01  11  10
 00  │ m0  m1  m3  m2                          00  │ 1   0   0   1      corners → B'D'
 01  │ m4  m5  m7  m6                          01  │ 0   1   1   0      middle  → BD
 11  │ m12 m13 m15 m14                         11  │ 0   1   1   0
 10  │ m8  m9  m11 m10                         10  │ 1   0   0   1      F = BD + B'D' = B ⊙ D
\end{Verbatim}
\end{listingblock}

Rules: group 1, 2, 4, 8, 16 adjacent cells (wrap-around allowed); larger groups = fewer literals; use don\textquotesingle t-cares (X) when helpful.

\begin{itemize}
\tightlist
\item
  \textbf{Prime implicant}: largest possible group. \textbf{Essential PI}: covers a minterm no other PI covers.
\end{itemize}

\needspace{114pt}

\hypertarget{p2-02--13-combinational-circuits}{%
\subsection{Combinational Circuits}\label{p2-02--13-combinational-circuits}}

\begin{listingblock}{64pt}
\begin{Verbatim}[fontsize=\fontsize{9.0}{10.9}\selectfont]
 Half adder                       Full adder
 S = A ⊕ B                        S    = A ⊕ B ⊕ Cin
 C = AB                           Cout = AB + Cin(A ⊕ B)
                                  (2 half adders + 1 OR gate)
\end{Verbatim}
\end{listingblock}

\begingroup\small

\begin{longtable}[]{@{}
  >{\raggedright\arraybackslash}p{(\columnwidth - 2\tabcolsep) * \real{0.2061}}
  >{\raggedright\arraybackslash}p{(\columnwidth - 2\tabcolsep) * \real{0.7939}}@{}}
\toprule\noalign{}
\begin{minipage}[b]{\linewidth}\raggedright
\textbf{Circuit}
\end{minipage} & \begin{minipage}[b]{\linewidth}\raggedright
\textbf{Function}
\end{minipage} \\
\midrule\noalign{}
\endhead
\bottomrule\noalign{}
\endlastfoot
Multiplexer 2ⁿ:1 & n select lines choose one of 2ⁿ inputs; any n-var function with 2ⁿ:1 MUX (or n+1 var with 2ⁿ:1 + inverter) \\
Demultiplexer 1:2ⁿ & Routes one input to one of 2ⁿ outputs \\
Decoder n:2ⁿ & Activates one of 2ⁿ outputs (minterm generator) \\
Encoder 2ⁿ:n & Reverse; priority encoder handles multiple inputs \\
Comparator & A\textgreater B, A=B, A\textless B \\
Ripple carry adder & n full adders; delay ∝ n \\
Carry look-ahead adder & Gᵢ = AᵢBᵢ, Pᵢ = Aᵢ⊕Bᵢ; Cᵢ₊₁ = Gᵢ + PᵢCᵢ --- faster \\
\end{longtable}

\endgroup

\needspace{130pt}

\hypertarget{p2-02--14-flip-flops}{%
\subsection{Flip-Flops}\label{p2-02--14-flip-flops}}

\begingroup\small

\begin{longtable}[]{@{}
  >{\raggedright\arraybackslash}p{(\columnwidth - 4\tabcolsep) * \real{0.0520}}
  >{\raggedright\arraybackslash}p{(\columnwidth - 4\tabcolsep) * \real{0.2425}}
  >{\raggedright\arraybackslash}p{(\columnwidth - 4\tabcolsep) * \real{0.4506}}@{}}
\toprule\noalign{}
\begin{minipage}[b]{\linewidth}\raggedright
\textbf{FF}
\end{minipage} & \begin{minipage}[b]{\linewidth}\raggedright
\textbf{Characteristic equation}
\end{minipage} & \begin{minipage}[b]{\linewidth}\raggedright
\textbf{Notes}
\end{minipage} \\
\midrule\noalign{}
\endhead
\bottomrule\noalign{}
\endlastfoot
SR & Q⁺ = S + R\textquotesingle Q (SR = 0) & S=R=1 invalid \\
JK & Q⁺ = JQ\textquotesingle{} + K\textquotesingle Q & J=K=1 toggles; race-around fixed by master-slave \\
D & Q⁺ = D & Delay / data latch \\
T & Q⁺ = T ⊕ Q & Toggle; used in counters \\
\end{longtable}

\endgroup

\begin{listingblock}{86pt}
\begin{Verbatim}[fontsize=\fontsize{9.0}{10.9}\selectfont]
 Excitation table (what inputs give Q → Q⁺)
 Q Q⁺ │ S R │ J K │ D │ T
 0 0  │ 0 X │ 0 X │ 0 │ 0
 0 1  │ 1 0 │ 1 X │ 1 │ 1
 1 0  │ 0 1 │ X 1 │ 0 │ 1
 1 1  │ X 0 │ X 0 │ 1 │ 0
\end{Verbatim}
\end{listingblock}

\needspace{95pt}

\hypertarget{p2-02--15-counters--registers}{%
\subsection{Counters \& Registers}\label{p2-02--15-counters--registers}}

\begin{itemize}
\tightlist
\item
  n flip-flops → \textbf{mod-2ⁿ} ripple counter. Ring counter with n FFs → mod-n. \textbf{Johnson (twisted ring)} with n FFs → \textbf{mod-2n}.
\item
  Asynchronous (ripple) counter delay = n × t\_pd; synchronous counters clock all FFs together.
\item
  Shift registers: SISO, SIPO, PISO, PIPO.
\item
  Memory: 2ᵏ × n RAM needs k address lines and n data lines. e.g. 4K × 8 → 12 address lines.
\item
  Number of 1K × 4 chips to build 16K × 8: (16 × 8)/(1 × 4) = \textbf{32}.
\end{itemize}

\needspace{166pt}

\hypertarget{p2-02--2-data-representation}{%
\section{Data Representation}\label{p2-02--2-data-representation}}

\needspace{130pt}

\hypertarget{p2-02--21-signed-numbers-n-bits}{%
\subsection{Signed numbers (n bits)}\label{p2-02--21-signed-numbers-n-bits}}

\begingroup\small

\begin{longtable}[]{@{}
  >{\raggedright\arraybackslash}p{(\columnwidth - 4\tabcolsep) * \real{0.1608}}
  >{\raggedright\arraybackslash}p{(\columnwidth - 4\tabcolsep) * \real{0.2713}}
  >{\raggedright\arraybackslash}p{(\columnwidth - 4\tabcolsep) * \real{0.1454}}@{}}
\toprule\noalign{}
\begin{minipage}[b]{\linewidth}\raggedright
\textbf{Representation}
\end{minipage} & \begin{minipage}[b]{\linewidth}\raggedright
\textbf{Range}
\end{minipage} & \begin{minipage}[b]{\linewidth}\raggedright
\textbf{Zeros}
\end{minipage} \\
\midrule\noalign{}
\endhead
\bottomrule\noalign{}
\endlastfoot
Sign-magnitude & −(2ⁿ⁻¹−1) \ldots{} +(2ⁿ⁻¹−1) & Two (+0, −0) \\
1\textquotesingle s complement & −(2ⁿ⁻¹−1) \ldots{} +(2ⁿ⁻¹−1) & Two \\
\textbf{2\textquotesingle s complement} & \textbf{−2ⁿ⁻¹ \ldots{} +(2ⁿ⁻¹−1)} & One \\
\end{longtable}

\endgroup

8-bit 2\textquotesingle s complement: −128 to +127. −5 = 11111011.
\textbf{Overflow} in 2\textquotesingle s-complement addition: carry into MSB ≠ carry out of MSB (or adding two same-sign numbers gives opposite sign).
r\textquotesingle s complement of N (n digits) = rⁿ − N; (r−1)\textquotesingle s complement = rⁿ − 1 − N.

\needspace{136pt}

\hypertarget{p2-02--22-floating-point--ieee-754}{%
\subsection{Floating Point --- IEEE 754}\label{p2-02--22-floating-point--ieee-754}}

\begin{listingblock}{86pt}
\begin{Verbatim}[fontsize=\fontsize{9.0}{10.9}\selectfont]
 Single precision (32 bits)
 ┌───┬──────────────┬────────────────────────────────┐
 │ S │ Exponent (8) │ Mantissa / fraction (23)       │   bias = 127
 └───┴──────────────┴────────────────────────────────┘
 Double precision (64 bits): 1 | 11 | 52            bias = 1023
 Value = (−1)^S × 1.M × 2^(E − bias)
\end{Verbatim}
\end{listingblock}

Example: −6.25 = −110.01₂ = −1.1001 × 2² → S=1, E = 129 = 10000001, M = 1001000\ldots{}
→ \textbf{C0C80000} (hex).

\begin{itemize}
\tightlist
\item
  E = 0, M = 0 → ±0; E = all 1s, M = 0 → ±∞; E = all 1s, M ≠ 0 → NaN; E = 0, M ≠ 0 → denormal.
\end{itemize}

\needspace{95pt}

\hypertarget{p2-02--23-codes}{%
\subsection{Codes}\label{p2-02--23-codes}}

\begin{itemize}
\tightlist
\item
  \textbf{BCD} (8421), \textbf{Excess-3} (BCD + 3, self-complementing), \textbf{Gray code} (one bit changes; G = B ⊕ (B \textgreater\textgreater{} 1)).
\item
  \textbf{Parity} detects single-bit (odd number of) errors.
\item
  \textbf{Hamming code}: number of parity bits r satisfies 2ʳ ≥ m + r + 1; parity bits at positions 1, 2, 4, 8 \ldots; detects 2, corrects 1 (with extra bit: SEC-DED).
\item
  \textbf{CRC}: polynomial division; used in networks.
\end{itemize}

\needspace{95pt}

\hypertarget{p2-02--24-arithmetic}{%
\subsection{Arithmetic}\label{p2-02--24-arithmetic}}

\begin{itemize}
\tightlist
\item
  \textbf{Booth\textquotesingle s algorithm} (signed multiplication): look at Qₙ Qₙ₊₁ → 10: subtract M; 01: add M; 00/11: shift only; arithmetic right shift each step.
\item
  Division: restoring and non-restoring.
\end{itemize}

\needspace{95pt}

\hypertarget{p2-02--3-register-transfer--micro-operations}{%
\section{Register Transfer \& Micro-operations}\label{p2-02--3-register-transfer--micro-operations}}

\begin{itemize}
\tightlist
\item
  RTL: \texttt{R2\ ←\ R1}, \texttt{P:\ R2\ ←\ R1} (conditional on control P).
\item
  Micro-operation types: register transfer, arithmetic (add, sub, inc, dec), logic, shift (logical, circular, arithmetic).
\item
  Common bus with k registers of n bits: n multiplexers of k×1 each, log₂k select lines.
\end{itemize}

\needspace{95pt}

\hypertarget{p2-02--4-basic-computer-mano}{%
\section{Basic Computer (Mano)}\label{p2-02--4-basic-computer-mano}}

\diagram{figures/p2-02-00}

Instruction format (16 bits): \texttt{I\ (1)\ \textbar{}\ Opcode\ (3)\ \textbar{}\ Address\ (12)} --- I = 0 direct, 1 indirect.

\begin{itemize}
\tightlist
\item
  Instruction types: memory-reference (opcode 000--110), register-reference (0111\ldots), I/O (1111\ldots).
\end{itemize}

\needspace{95pt}

\hypertarget{p2-02--instruction-cycle}{%
\subsection{Instruction Cycle}\label{p2-02--instruction-cycle}}

\diagram{figures/p2-02-01}

\needspace{130pt}

\hypertarget{p2-02--5-control-unit}{%
\section{Control Unit}\label{p2-02--5-control-unit}}

\begingroup\small

\begin{longtable}[]{@{}
  >{\raggedright\arraybackslash}p{(\columnwidth - 2\tabcolsep) * \real{0.3500}}
  >{\raggedright\arraybackslash}p{(\columnwidth - 2\tabcolsep) * \real{0.3594}}@{}}
\toprule\noalign{}
\begin{minipage}[b]{\linewidth}\raggedright
\textbf{Hardwired}
\end{minipage} & \begin{minipage}[b]{\linewidth}\raggedright
\textbf{Microprogrammed}
\end{minipage} \\
\midrule\noalign{}
\endhead
\bottomrule\noalign{}
\endlastfoot
Logic gates, decoders, sequence counter & Control memory stores microinstructions \\
Fast & Slower \\
Difficult to modify & Flexible, easy to modify \\
Used in RISC & Used in CISC \\
\end{longtable}

\endgroup

\begin{itemize}
\tightlist
\item
  Control word bits: \textbf{horizontal} microprogramming (one bit per control signal, wide, parallel, fast) vs \textbf{vertical} (encoded, narrow, needs decoder).
\item
  Control address register (CAR), control data register (CDR), sequencer, mapping logic.
\item
  Control memory size = (\#microinstructions) × (width).
\end{itemize}

\needspace{196pt}

\hypertarget{p2-02--6-cpu-organisation}{%
\section{CPU Organisation}\label{p2-02--6-cpu-organisation}}

\needspace{160pt}

\hypertarget{p2-02--61-instruction-formats-by-number-of-addresses}{%
\subsection{Instruction formats (by number of addresses)}\label{p2-02--61-instruction-formats-by-number-of-addresses}}

Evaluate X = (A + B) × (C + D):

\begingroup\small

\begin{longtable}[]{@{}
  >{\raggedright\arraybackslash}p{(\columnwidth - 2\tabcolsep) * \real{0.2206}}
  >{\raggedright\arraybackslash}p{(\columnwidth - 2\tabcolsep) * \real{0.7794}}@{}}
\toprule\noalign{}
\begin{minipage}[b]{\linewidth}\raggedright
\textbf{Type}
\end{minipage} & \begin{minipage}[b]{\linewidth}\raggedright
\textbf{Code}
\end{minipage} \\
\midrule\noalign{}
\endhead
\bottomrule\noalign{}
\endlastfoot
3-address & ADD R1, A, B · ADD R2, C, D · MUL X, R1, R2 \\
2-address & MOV R1, A · ADD R1, B · MOV R2, C · ADD R2, D · MUL R1, R2 · MOV X, R1 \\
1-address (accumulator) & LOAD A · ADD B · STORE T · LOAD C · ADD D · MUL T · STORE X \\
0-address (stack) & PUSH A · PUSH B · ADD · PUSH C · PUSH D · ADD · MUL · POP X \\
\end{longtable}

\endgroup

\needspace{130pt}

\hypertarget{p2-02--62-addressing-modes-very-frequent}{%
\subsection{Addressing Modes (very frequent)}\label{p2-02--62-addressing-modes-very-frequent}}

\begingroup\small

\begin{longtable}[]{@{}
  >{\raggedright\arraybackslash}p{(\columnwidth - 4\tabcolsep) * \real{0.2477}}
  >{\raggedright\arraybackslash}p{(\columnwidth - 4\tabcolsep) * \real{0.2080}}
  >{\raggedright\arraybackslash}p{(\columnwidth - 4\tabcolsep) * \real{0.3293}}@{}}
\toprule\noalign{}
\begin{minipage}[b]{\linewidth}\raggedright
\textbf{Mode}
\end{minipage} & \begin{minipage}[b]{\linewidth}\raggedright
\textbf{Effective address}
\end{minipage} & \begin{minipage}[b]{\linewidth}\raggedright
\textbf{Use}
\end{minipage} \\
\midrule\noalign{}
\endhead
\bottomrule\noalign{}
\endlastfoot
Implied & Operand implicit & CMA, stack ops \\
Immediate & Operand in instruction & Constants \\
Register & Operand in register & Fast \\
Register indirect & EA = (R) & Pointers \\
Direct (absolute) & EA = A & Global variables \\
Indirect & EA = M{[}A{]} & Pointers \\
Auto-increment/\allowbreak{}decrement & EA = (R), then R±1 & Arrays, stacks \\
Relative & EA = PC + A & Branches, \textbf{position-independent code} \\
Indexed & EA = XR + A & Arrays \\
Base register & EA = BR + A & Relocation \\
\end{longtable}

\endgroup

Memory references to get operand: immediate 0, direct 1, indirect 2.

\needspace{130pt}

\hypertarget{p2-02--63-risc-vs-cisc}{%
\subsection{RISC vs CISC}\label{p2-02--63-risc-vs-cisc}}

\begingroup\small

\begin{longtable}[]{@{}
  >{\raggedright\arraybackslash}p{(\columnwidth - 2\tabcolsep) * \real{0.3472}}
  >{\raggedright\arraybackslash}p{(\columnwidth - 2\tabcolsep) * \real{0.3317}}@{}}
\toprule\noalign{}
\begin{minipage}[b]{\linewidth}\raggedright
\textbf{RISC}
\end{minipage} & \begin{minipage}[b]{\linewidth}\raggedright
\textbf{CISC}
\end{minipage} \\
\midrule\noalign{}
\endhead
\bottomrule\noalign{}
\endlastfoot
Few, simple, fixed-length instructions & Many, complex, variable length \\
Few addressing modes & Many addressing modes \\
Load/store architecture & Memory operands in many instructions \\
Many registers, register windows & Fewer registers \\
Hardwired control & Microprogrammed \\
1 instruction per cycle, pipelining easy & Multi-cycle \\
ARM, MIPS, SPARC, RISC-V & x86, VAX, IBM 370 \\
\end{longtable}

\endgroup

\needspace{158pt}

\hypertarget{p2-02--7-pipelining}{%
\section{Pipelining}\label{p2-02--7-pipelining}}

\begin{listingblock}{108pt}
\begin{Verbatim}[fontsize=\fontsize{9.0}{10.9}\selectfont]
 4-stage instruction pipeline (IF, ID, EX, WB), 5 instructions
 Cycle:   1    2    3    4    5    6    7    8
 I1      IF   ID   EX   WB
 I2           IF   ID   EX   WB
 I3                IF   ID   EX   WB
 I4                     IF   ID   EX   WB
 I5                          IF   ID   EX   WB
 Total = k + (n − 1) = 4 + 4 = 8 cycles
\end{Verbatim}
\end{listingblock}

\begin{listingblock}{79pt}
\begin{Verbatim}[fontsize=\fontsize{8.1}{9.8}\selectfont]
Pipelined time   T_pipe = (k + n − 1) × t_p        (t_p = max stage delay + latch delay)
Non-pipelined    T_seq  = n × t_n                   (t_n = sum of stage delays)
Speedup          S = n·t_n / ((k + n − 1)·t_p)  → k  as n → ∞ (ideal, equal stages)
Efficiency       η = S / k
Throughput       n / T_pipe
With stalls:     CPI_pipe = 1 + stall cycles per instruction ;  S = CPI_nonpipe / CPI_pipe
\end{Verbatim}
\end{listingblock}

\needspace{95pt}

\hypertarget{p2-02--hazards}{%
\subsection{Hazards}\label{p2-02--hazards}}

\diagram{figures/p2-02-02}

\begin{itemize}
\tightlist
\item
  Only \textbf{RAW} occurs in a simple in-order pipeline; WAR and WAW occur in out-of-order execution.
\item
  \textbf{Flynn\textquotesingle s taxonomy}: SISD (uniprocessor), \textbf{SIMD} (vector/array processors, GPU), MISD (rare; fault-tolerant), MIMD (multiprocessors).
\item
  \textbf{Vector processing}: one instruction operates on arrays; \textbf{array processors} = SIMD with multiple PEs.
\item
  \textbf{Amdahl\textquotesingle s law}: Speedup = 1 / ((1 − f) + f/s), f = fraction enhanced, s = speedup of that fraction.
\item
  Superscalar (multiple issue), VLIW, superpipelining.
\end{itemize}

\needspace{95pt}

\hypertarget{p2-02--8-io-organisation}{%
\section{I/O Organisation}\label{p2-02--8-io-organisation}}

\diagram{figures/p2-02-03}

\begingroup\small

\begin{longtable}[]{@{}
  >{\raggedright\arraybackslash}p{(\columnwidth - 4\tabcolsep) * \real{0.2305}}
  >{\raggedright\arraybackslash}p{(\columnwidth - 4\tabcolsep) * \real{0.4557}}
  >{\raggedright\arraybackslash}p{(\columnwidth - 4\tabcolsep) * \real{0.2339}}@{}}
\toprule\noalign{}
\begin{minipage}[b]{\linewidth}\raggedright
\textbf{Mode}
\end{minipage} & \begin{minipage}[b]{\linewidth}\raggedright
\textbf{How}
\end{minipage} & \begin{minipage}[b]{\linewidth}\raggedright
\textbf{CPU involvement}
\end{minipage} \\
\midrule\noalign{}
\endhead
\bottomrule\noalign{}
\endlastfoot
Programmed I/O & CPU polls status flag & Very high (busy waiting) \\
Interrupt-driven & Device interrupts CPU & Per word \\
\textbf{DMA} & DMA controller transfers block directly to memory & Only at start/end \\
I/O processor (channel) & Dedicated processor & Minimal \\
\end{longtable}

\endgroup

\begin{itemize}
\tightlist
\item
  DMA modes: \textbf{burst} (block at once), \textbf{cycle stealing} (one word per cycle), \textbf{transparent} (when CPU not using bus).
\item
  DMA signals: Bus Request (BR) and Bus Grant (BG).
\item
  Asynchronous transfer: \textbf{strobe} (one control line) or \textbf{handshaking} (two lines --- data valid \& data accepted).
\item
  \textbf{Priority interrupts}: polling (software), \textbf{daisy chaining} (serial hardware --- nearest device highest priority), parallel priority (priority encoder).
\item
  Vectored interrupt: device supplies ISR address. Non-maskable interrupt (NMI): cannot be disabled.
\item
  Isolated I/O (separate address space, IN/OUT instructions) vs \textbf{memory-mapped I/O} (same address space).
\item
  Serial communication: asynchronous (start bit, data, parity, stop bits) --- UART; synchronous --- USRT.
\end{itemize}

\needspace{190pt}

\hypertarget{p2-02--9-memory-hierarchy}{%
\section{Memory Hierarchy}\label{p2-02--9-memory-hierarchy}}

\begin{listingblock}{140pt}
\begin{Verbatim}[fontsize=\fontsize{9.0}{10.9}\selectfont]
 ┌─────────────────────────────┐  ▲  faster, costlier, smaller
 │ Registers                   │  │
 ├─────────────────────────────┤  │
 │ Cache L1 / L2 / L3 (SRAM)   │  │
 ├─────────────────────────────┤  │
 │ Main memory (DRAM)          │  │
 ├─────────────────────────────┤  │
 │ SSD / Magnetic disk         │  │
 ├─────────────────────────────┤  │
 │ Tape / Optical              │  ▼  slower, cheaper, larger
 └─────────────────────────────┘
\end{Verbatim}
\end{listingblock}

\begin{itemize}
\tightlist
\item
  SRAM (flip-flops, no refresh, cache) vs DRAM (capacitors, refresh, main memory).
\item
  ROM, PROM, EPROM (UV erase), EEPROM (electrical), Flash.
\item
  \textbf{Associative memory (CAM)}: accessed by content; used in TLB.
\item
  \textbf{Locality of reference}: temporal (same item soon) \& spatial (nearby items soon).
\end{itemize}

\needspace{95pt}

\hypertarget{p2-02--91-cache-mapping}{%
\subsection{Cache Mapping}\label{p2-02--91-cache-mapping}}

\diagram{figures/p2-02-04}

\textbf{Worked example}: Main memory 4 GB (32-bit address), cache 64 KB, block 32 B.

\begin{listingblock}{86pt}
\begin{Verbatim}[fontsize=\fontsize{9.0}{10.9}\selectfont]
Offset bits  = log₂ 32 = 5
Lines        = 64 KB / 32 B = 2048 → 11 bits
Direct:  TAG = 32 − 11 − 5 = 16 | LINE 11 | OFFSET 5
4-way:   sets = 2048/4 = 512 → 9 bits ; TAG = 32 − 9 − 5 = 18
Fully:   TAG = 32 − 5 = 27
Tag directory size (direct) = 2048 × 16 bits (+ valid/dirty bits)
\end{Verbatim}
\end{listingblock}

\begin{listingblock}{64pt}
\begin{Verbatim}[fontsize=\fontsize{9.0}{10.9}\selectfont]
Hit ratio h ;  cache access tc ; memory access tm
Avg access time (simultaneous/parallel)   T = h·tc + (1 − h)·tm
Avg access time (hierarchical/serial)     T = h·tc + (1 − h)(tc + tm)
AMAT = Hit time + Miss rate × Miss penalty
\end{Verbatim}
\end{listingblock}

\begin{itemize}
\tightlist
\item
  Replacement: LRU, FIFO, random (needed for associative/set-associative only).
\item
  Write policies: \textbf{write-through} (update memory immediately; with write buffer) vs \textbf{write-back} (dirty bit, on eviction).
\item
  Write miss: write-allocate (usually with write-back) vs no-write-allocate (with write-through).
\item
  Misses: \textbf{Compulsory} (cold), \textbf{Capacity}, \textbf{Conflict} (3 Cs).
\end{itemize}

\needspace{125pt}

\hypertarget{p2-02--92-virtual-memory}{%
\subsection{Virtual Memory}\label{p2-02--92-virtual-memory}}

Covered in depth in \protect\hyperlink{p2-05}{Unit 5 (OS)}: paging, TLB, page replacement.

\begin{itemize}
\tightlist
\item
  Effective access with TLB: EAT = h(t\_TLB + t\_m) + (1 − h)(t\_TLB + 2t\_m) (single-level page table).
\end{itemize}

\needspace{103pt}

\hypertarget{p2-02--93-magnetic-disk}{%
\subsection{Magnetic Disk}\label{p2-02--93-magnetic-disk}}

\begin{listingblock}{53pt}
\begin{Verbatim}[fontsize=\fontsize{9.0}{10.9}\selectfont]
Disk access time = Seek time + Rotational latency + Transfer time
Avg rotational latency = ½ × (60/RPM) s
Capacity = surfaces × tracks × sectors/track × bytes/sector
\end{Verbatim}
\end{listingblock}

e.g. 7200 RPM → one rotation 8.33 ms → avg latency \textbf{4.17 ms}.

\needspace{95pt}

\hypertarget{p2-02--10-multiprocessors}{%
\section{Multiprocessors}\label{p2-02--10-multiprocessors}}

\begin{itemize}
\tightlist
\item
  \textbf{Tightly coupled} (shared memory, UMA/NUMA) vs \textbf{loosely coupled} (distributed memory, message passing).
\item
  Interconnection structures: time-shared common bus, multiport memory, \textbf{crossbar switch} (n² switches, non-blocking), multistage networks (Omega: (n/2) log₂ n 2×2 switches), hypercube (n-cube: 2ⁿ nodes, each with n links, diameter n).
\item
  Inter-processor arbitration: serial (daisy chain), parallel, dynamic (time-slice, polling, LRU, FIFO, rotating daisy chain).
\item
  Synchronisation: test-and-set, semaphores, mutual exclusion via hardware lock.
\item
  \textbf{Cache coherence}: write-invalidate / write-update; snoopy protocols; \textbf{MESI} (Modified, Exclusive, Shared, Invalid); directory-based protocols.
\item
  \textbf{Multicore}: multiple cores on one chip sharing L2/L3 cache; hyper-threading (SMT).
\end{itemize}

\needspace{196pt}

\hypertarget{p2-02--11-deeper-dive--booth-trace-instruction-encoding-cache-behaviour--interrupts}{%
\section{Deeper Dive --- Booth Trace, Instruction Encoding, Cache Behaviour \& Interrupts}\label{p2-02--11-deeper-dive--booth-trace-instruction-encoding-cache-behaviour--interrupts}}

\needspace{160pt}

\hypertarget{p2-02--111-booths-algorithm--full-trace-7--3-4-bits}{%
\subsection{Booth\textquotesingle s Algorithm --- Full Trace: 7 × (−3), 4 bits}\label{p2-02--111-booths-algorithm--full-trace-7--3-4-bits}}

M = 0111 (7), −M = 1001, Q = 1101 (−3), A = 0000, Q₋₁ = 0.

\begingroup\small

\begin{longtable}[]{@{}
  >{\raggedright\arraybackslash}p{(\columnwidth - 10\tabcolsep) * \real{0.0709}}
  >{\raggedright\arraybackslash}p{(\columnwidth - 10\tabcolsep) * \real{0.1219}}
  >{\raggedright\arraybackslash}p{(\columnwidth - 10\tabcolsep) * \real{0.2321}}
  >{\raggedright\arraybackslash}p{(\columnwidth - 10\tabcolsep) * \real{0.0705}}
  >{\raggedright\arraybackslash}p{(\columnwidth - 10\tabcolsep) * \real{0.0705}}
  >{\raggedright\arraybackslash}p{(\columnwidth - 10\tabcolsep) * \real{0.0759}}@{}}
\toprule\noalign{}
\begin{minipage}[b]{\linewidth}\raggedright
\textbf{Step}
\end{minipage} & \begin{minipage}[b]{\linewidth}\raggedright
\textbf{Q₀ Q₋₁}
\end{minipage} & \begin{minipage}[b]{\linewidth}\raggedright
\textbf{Operation}
\end{minipage} & \begin{minipage}[b]{\linewidth}\raggedright
\textbf{A}
\end{minipage} & \begin{minipage}[b]{\linewidth}\raggedright
\textbf{Q}
\end{minipage} & \begin{minipage}[b]{\linewidth}\raggedright
\textbf{Q₋₁}
\end{minipage} \\
\midrule\noalign{}
\endhead
\bottomrule\noalign{}
\endlastfoot
Init & & & 0000 & 1101 & 0 \\
1 & 1 0 & A ← A − M & 1001 & 1101 & 0 \\
& & Arithmetic shift right & 1100 & 1110 & 1 \\
2 & 0 1 & A ← A + M & 0011 & 1110 & 1 \\
& & Arithmetic shift right & 0001 & 1111 & 0 \\
3 & 1 0 & A ← A − M & 1010 & 1111 & 0 \\
& & Arithmetic shift right & 1101 & 0111 & 1 \\
4 & 1 1 & Shift only & 1110 & 1011 & 1 \\
\end{longtable}

\endgroup

Result A Q = \textbf{1110 1011} = −21 in 8-bit 2\textquotesingle s complement ✔.
Booth reduces additions for runs of 1s; worst case is alternating bits (0101\ldots).

\needspace{122pt}

\hypertarget{p2-02--112-instruction-encoding-numericals}{%
\subsection{Instruction Encoding Numericals}\label{p2-02--112-instruction-encoding-numericals}}

\textbf{Fixed format}: 32-bit instructions, 64 registers, 45 distinct opcodes, format \texttt{opcode\ \textbar{}\ Rd\ \textbar{}\ Rs\ \textbar{}\ immediate}.

\begin{listingblock}{42pt}
\begin{Verbatim}[fontsize=\fontsize{9.0}{10.9}\selectfont]
opcode bits = ⌈log₂ 45⌉ = 6        register field = log₂ 64 = 6 each
immediate   = 32 − 6 − 6 − 6 = 14 bits  → signed range −8192 … +8191
\end{Verbatim}
\end{listingblock}

\textbf{Expanding opcode}: 16-bit instruction, 4-bit address fields.

\begin{listingblock}{68pt}
\begin{Verbatim}[fontsize=\fontsize{7.9}{9.6}\selectfont]
3-address: 4-bit opcode → 2⁴ = 16 patterns; use 15, keep 1 as escape
2-address: escape (4 bits) + 4 more opcode bits → 16 patterns; use 14, keep 2
1-address: 2 escapes × 16 = 32 patterns; use 31, keep 1
0-address: 1 escape × 16 = 16 instructions
General rule: unused patterns at level k × 2^(field width) = patterns available at level k+1
\end{Verbatim}
\end{listingblock}

\needspace{133pt}

\hypertarget{p2-02--113-mapping-an-address-into-a-set-associative-cache}{%
\subsection{Mapping an Address into a Set-Associative Cache}\label{p2-02--113-mapping-an-address-into-a-set-associative-cache}}

2-way set associative, 128 lines, 16-byte blocks, 16-bit byte address \textbf{0x1A2B}.

\begin{listingblock}{53pt}
\begin{Verbatim}[fontsize=\fontsize{9.0}{10.9}\selectfont]
Sets = 128 / 2 = 64 → 6 set bits ; offset = 4 bits ; tag = 16 − 6 − 4 = 6 bits
Block number = 0x1A2B >> 4 = 0x1A2 = 418
Set   = 418 mod 64 = 34           Tag = 418 div 64 = 6         Offset = 0xB = 11
\end{Verbatim}
\end{listingblock}

\needspace{160pt}

\hypertarget{p2-02--114-hits--misses-for-three-organisations}{%
\subsection{Hits \& Misses for Three Organisations}\label{p2-02--114-hits--misses-for-three-organisations}}

Block reference string: \textbf{0, 4, 0, 4, 8, 0} with 4 cache lines.

\begingroup\small

\begin{longtable}[]{@{}
  >{\raggedright\arraybackslash}p{(\columnwidth - 4\tabcolsep) * \real{0.4034}}
  >{\raggedright\arraybackslash}p{(\columnwidth - 4\tabcolsep) * \real{0.5124}}
  >{\raggedright\arraybackslash}p{(\columnwidth - 4\tabcolsep) * \real{0.0842}}@{}}
\toprule\noalign{}
\begin{minipage}[b]{\linewidth}\raggedright
\textbf{Organisation}
\end{minipage} & \begin{minipage}[b]{\linewidth}\raggedright
\textbf{Trace}
\end{minipage} & \begin{minipage}[b]{\linewidth}\raggedright
\textbf{Misses}
\end{minipage} \\
\midrule\noalign{}
\endhead
\bottomrule\noalign{}
\endlastfoot
Direct mapped (line = block mod 4) & 0, 4, 8 all map to line 0 → every access evicts the previous & \textbf{6} \\
2-way set assoc., LRU (set = block mod 2) & 0 M, 4 M, 0 H, 4 H, 8 M (evict 0), 0 M (evict 4) & \textbf{4} \\
Fully associative, LRU & 0 M, 4 M, 0 H, 4 H, 8 M, 0 H & \textbf{3} \\
\end{longtable}

\endgroup

Higher associativity removes \textbf{conflict misses}; the three first-time misses are \textbf{compulsory}.

\needspace{95pt}

\hypertarget{p2-02--115-interrupt-cycle-mano-basic-computer}{%
\subsection{Interrupt Cycle (Mano basic computer)}\label{p2-02--115-interrupt-cycle-mano-basic-computer}}

\diagram{figures/p2-02-05}

Return address is saved in memory location 0; the ISR ends with an indirect branch through location 0 (BUN 0 I), re-enabling interrupts with ION.

\needspace{125pt}

\hypertarget{p2-02--116-microprogrammed-control-numericals}{%
\subsection{Microprogrammed Control Numericals}\label{p2-02--116-microprogrammed-control-numericals}}

Mano\textquotesingle s microinstruction: \texttt{F1\ (3)\ \textbar{}\ F2\ (3)\ \textbar{}\ F3\ (3)\ \textbar{}\ CD\ (2)\ \textbar{}\ BR\ (2)\ \textbar{}\ AD\ (7)} = \textbf{20 bits}; control memory 128 × 20.

\begin{itemize}
\tightlist
\item
  Horizontal: n control signals → n bits. Vertical: encode k mutually exclusive signals in ⌈log₂(k + 1)⌉ bits.
\item
  \textbf{Example}: 48 control signals in 4 mutually-exclusive groups of 12 → vertical needs 4 × ⌈log₂ 13⌉ = 4 × 4 = 16 bits (vs 48 horizontal).
\end{itemize}

\needspace{95pt}

\hypertarget{p2-02--117-dma--io-numericals}{%
\subsection{DMA \& I/O Numericals}\label{p2-02--117-dma--io-numericals}}

\begin{itemize}
\tightlist
\item
  Disk transfers 4 MB/s in 32-bit words using cycle stealing on a memory that supports 100 million cycles/s:
  words/s = 4 MB / 4 B = 1 M → fraction of memory cycles stolen = 1 M / 100 M = \textbf{1\%}.
\item
  Programmed I/O polling overhead: polling 100 times/s × 500 cycles each = 50,000 cycles/s → \textbf{0.05\%} of a 100 MHz CPU.
\end{itemize}

\needspace{123pt}

\hypertarget{p2-02--118-memory-interleaving}{%
\subsection{Memory Interleaving}\label{p2-02--118-memory-interleaving}}

\begin{listingblock}{73pt}
\begin{Verbatim}[fontsize=\fontsize{8.7}{10.6}\selectfont]
 Low-order interleaving (4 modules): module = address mod 4
 Address:  0  1  2  3  4  5  6  7  8 …
 Module:   M0 M1 M2 M3 M0 M1 M2 M3 M0 …
 Consecutive words come from different modules → overlapped access, higher bandwidth
 High-order interleaving: module = high bits → consecutive words in the same module
\end{Verbatim}
\end{listingblock}

\needspace{125pt}

\hypertarget{p2-02--previous-year-questions-pyq-pattern}{%
\section{Previous Year Questions (PYQ pattern)}\label{p2-02--previous-year-questions-pyq-pattern}}

\textbf{Number systems \& codes}

\begin{enumerate}
\def\labelenumi{\arabic{enumi}.}
\tightlist
\item
  2\textquotesingle s complement of 01101000: \textbf{10011000}
\item
  Range of 8-bit 2\textquotesingle s complement: \textbf{−128 to +127}
\item
  Gray code of binary 1011: 1, 1⊕0, 0⊕1, 1⊕1 → \textbf{1110}
\item
  Excess-3 of decimal 5: \textbf{1000}
\item
  Hamming code for 4 data bits needs \textbf{3} parity bits (2³ ≥ 4+3+1)
\item
  IEEE 754 single precision bias: \textbf{127}; exponent bits \textbf{8}; mantissa \textbf{23}
\item
  (−0.75) in IEEE 754 single: S=1, −1.1×2⁻¹, E=126 → \textbf{BF400000}
\end{enumerate}

\textbf{Digital logic}

\begin{enumerate}
\def\labelenumi{\arabic{enumi}.}
\setcounter{enumi}{7}
\tightlist
\item
  Minimum number of NAND gates for XOR: \textbf{4}
\item
  Simplify F = Σm(0,1,2,3): \textbf{A\textquotesingle{}} (for 3 variables A,B,C) --- F = A\textquotesingle{}
\item
  Number of select lines for 32:1 MUX: \textbf{5}
\item
  Johnson counter with 4 FFs has \textbf{8} states
\item
  Characteristic equation of JK FF: \textbf{Q⁺ = JQ\textquotesingle{} + K\textquotesingle Q}
\item
  Race-around condition occurs in: \textbf{JK flip-flop when J=K=1 and clock pulse wide} --- solved by master-slave
\item
  A 3-to-8 decoder with an OR gate can implement: \textbf{any 3-variable Boolean function}
\item
  16K × 8 memory needs \textbf{14} address lines
\end{enumerate}

\textbf{CPU \& addressing}

\begin{enumerate}
\def\labelenumi{\arabic{enumi}.}
\setcounter{enumi}{15}
\tightlist
\item
  Addressing mode used for position-independent code: \textbf{Relative (PC-relative)}
\item
  Addressing mode in which operand is part of instruction: \textbf{Immediate}
\item
  Which addressing mode needs two memory accesses to get operand (after fetch)? \textbf{Indirect}
\item
  Zero-address instructions are used in: \textbf{stack-organised computers}
\item
  Which is a feature of RISC? \textbf{Fixed-length instructions / load-store / hardwired control}
\item
  Horizontal microprogramming: \textbf{one bit per control signal, longer control words}
\end{enumerate}

\textbf{Pipelining}

\begin{enumerate}
\def\labelenumi{\arabic{enumi}.}
\setcounter{enumi}{21}
\tightlist
\item
  A 5-stage pipeline, stage delay 20 ns, executes 100 instructions. Time = (5 + 99) × 20 = \textbf{2080 ns}
\item
  Ideal speedup of a k-stage pipeline: \textbf{k}
\item
  Stages 150, 120, 160, 140 ns, latch 5 ns. Pipeline cycle = \textbf{165 ns}; non-pipelined instruction = 570 ns
\item
  RAW hazard is also called: \textbf{true data dependency}; fixed by \textbf{operand forwarding}
\item
  Amdahl: 40\% of program sped up 10×. Speedup = 1/(0.6 + 0.04) = \textbf{1.5625}
\item
  Vector processors belong to which Flynn category? \textbf{SIMD}
\end{enumerate}

\textbf{Memory \& cache}

\begin{enumerate}
\def\labelenumi{\arabic{enumi}.}
\setcounter{enumi}{27}
\tightlist
\item
  Cache 2 ns, memory 20 ns, hit ratio 0.9, hierarchical: 0.9×2 + 0.1×22 = \textbf{4 ns}
\item
  Direct-mapped cache with 128 lines, block 16 words; 64K-word memory → address 16 bits: TAG \textbf{5}, LINE \textbf{7}, WORD \textbf{4}
\item
  Which miss cannot be reduced by increasing cache size? \textbf{Compulsory miss}
\item
  Write-back cache uses a: \textbf{dirty bit}
\item
  Associative memory is accessed by: \textbf{content}
\item
  Disk 6000 RPM: average rotational latency = \textbf{5 ms}
\end{enumerate}

\textbf{I/O}

\begin{enumerate}
\def\labelenumi{\arabic{enumi}.}
\setcounter{enumi}{33}
\tightlist
\item
  Which mode transfers data without CPU intervention? \textbf{DMA}
\item
  In daisy chaining, priority is determined by: \textbf{position of device in the chain}
\item
  DMA cycle stealing means: \textbf{DMA takes the bus for one cycle at a time}
\item
  Handshaking uses: \textbf{two control lines}
\end{enumerate}

\textbf{Multiprocessors}

\begin{enumerate}
\def\labelenumi{\arabic{enumi}.}
\setcounter{enumi}{37}
\tightlist
\item
  Number of crosspoints in a crossbar connecting 8 processors to 8 memories: \textbf{64}
\item
  MESI protocol is used for: \textbf{cache coherence}
\item
  In a 4-dimensional hypercube, number of nodes: \textbf{16}, links per node: \textbf{4}
\end{enumerate}

\textbf{More practice questions}

\begin{enumerate}
\def\labelenumi{\arabic{enumi}.}
\setcounter{enumi}{40}
\tightlist
\item
  Booth\textquotesingle s algorithm when Q₀Q₋₁ = 01: \textbf{add M then shift}
\item
  A machine has 32-bit instructions and 128 registers; a 3-register format leaves how many opcode bits? 32 − 21 = \textbf{11}
\item
  16-bit instructions with 6-bit address fields: 2-address instructions use 4-bit opcodes; with 14 two-address instructions used, one-address instructions possible: 2 × 64 = \textbf{128}
\item
  4-way set associative cache, 256 lines, 32-byte blocks, 32-bit addresses: set bits \textbf{6}, offset \textbf{5}, tag \textbf{21}
\item
  Which miss is eliminated by full associativity? \textbf{conflict miss}
\item
  In Mano\textquotesingle s interrupt cycle, the return address is stored in: \textbf{memory location 0}
\item
  Vertical microprogramming requires a: \textbf{decoder} for the encoded control fields
\item
  Low-order interleaving places consecutive addresses in: \textbf{different memory modules}
\item
  A DMA transfer of 1 M words/s on 50 M memory cycles/s steals: \textbf{2\%} of cycles
\item
  Booth multiplication of 2 × (−4) in 4 bits gives A Q = \textbf{1111 1000} (−8)
\end{enumerate}

\begin{revbox}

\begin{itemize}
\tightlist
\item
  NAND: XOR=4 · NOR: XNOR=4 · Johnson n FF → 2n states
\item
  IEEE single 1/8/23 bias 127 · double 1/11/52 bias 1023
\item
  Pipeline time (k + n − 1)t · speedup → k · Amdahl 1/((1−f)+f/s)
\item
  Hierarchical cache T = h·tc + (1−h)(tc + tm)
\item
  Direct: TAG\textbar LINE\textbar OFFSET · Set-assoc: TAG\textbar SET\textbar OFFSET · Fully: TAG\textbar OFFSET
\end{itemize}

\end{revbox}
